% ============================================================
\chapter{Código fuente y datos analizados}
\label{ap:repositorio}
% ============================================================

El código fuente completo del sistema \textit{kubescan}, incluyendo los
\textit{scripts} de adquisición de datos, entrenamiento de modelos,
construcción de grafos y la interfaz de línea de comandos, se encuentra
disponible en el siguiente repositorio de GitHub:

\begin{center}
  \url{https://github.com/ObedRav/kubescan}
\end{center}

El repositorio incluye los siguientes componentes:

\begin{itemize}
  \item \textbf{\texttt{kubescan/}}: paquete Python instalable con la \ac{CLI} y
    los modelos de inferencia.
  \item \textbf{\texttt{kubescan/tests/}}: batería de pruebas automatizadas, con
    siete módulos unitarios (\texttt{test\_yaml\_parser.py},
    \texttt{test\_graph\_builder.py}, \texttt{test\_gat\_encoder.py},
    \texttt{test\_rf\_classifier.py}, \texttt{test\_ga\_ensemble.py},
    \texttt{test\_device\_utils.py} y \texttt{test\_properties.py}, este último
    basado en pruebas de propiedades con \textit{Hypothesis}) y una prueba de
    integración de extremo a extremo sobre la \ac{CLI}
    (\texttt{integration/test\_cli\_scan.py}).
  \item \textbf{\texttt{research/scripts/}}: \textit{scripts} de adquisición y
    preprocesamiento de los manifiestos \ac{YAML} (fases 1 y 2).
  \item \textbf{\texttt{research/models/}}: código de entrenamiento del Random
    Forest (\texttt{train\_rf.py}), de la \ac{GAT} (\texttt{train\_gnn.py}), y del
    Algoritmo Genético (\texttt{run\_ga\_ensemble.py}), junto a los \textit{scripts}
    de verificación que reproducen las comprobaciones citadas en el
    Capítulo~\ref{cap:conclusiones}: \texttt{verify\_label\_circularity.py}, que
    contrasta la etiqueta de la Capa~1 con la disyunción de las columnas de
    misconfiguración, y \texttt{verify\_risk\_score\_channel.py}, que documenta la
    procedencia y la cobertura del atributo nodal~25 y ejecuta su ablación sobre
    el conjunto de test.
  \item \textbf{\texttt{research/data/}}: los grafos de clúster en formato
    \texttt{.npz} (NumPy, con campos \texttt{x}, \texttt{edge\_index},
    \texttt{edge\_attr}, \texttt{y}, \texttt{node\_labels} y \texttt{risk\_scores})
    y el dataset tabular en formato \texttt{.csv}.
  \item \textbf{\texttt{research/models/checkpoints/}}: \textit{checkpoints}
    de los 5~modelos \ac{GAT} entrenados en validación cruzada y el modelo
    \ac{RF} (disponible también como enlace simbólico en
    \texttt{kubescan/checkpoints/trained/}).
  \item \textbf{\texttt{thesis/usability-study/}}: protocolo, cuestionario e
    instrumentos del estudio de usabilidad descrito en la
    Sección~\ref{sec:usabilidad}, junto con la plantilla de resultados y el
    \textit{script} de agregación \texttt{compute\_results.py}.
  \item \textbf{\texttt{Makefile}}: objetivos de reproducción del \textit{pipeline}
    de investigación (\texttt{make data}, \texttt{make reproduce}).
  \item \textbf{\texttt{requirements-lock.txt}}: fijación exacta de versiones de
    todas las dependencias transitivas del entorno de entrenamiento.
  \item \textbf{\texttt{.pre-commit-config.yaml}} y
    \textbf{\texttt{.github/workflows/ci.yml}}: controles automáticos de estilo,
    tipado y ejecución de la batería de pruebas en integración continua.
  \item \textbf{\texttt{thesis/}}: fuentes LaTeX de esta memoria.
\end{itemize}

Los pesos entrenados se distribuyen junto con el código, de modo que un clon
nuevo del repositorio puede ejecutar \texttt{kubescan scan} sin reentrenar
ningún modelo. El directorio \texttt{research/models/checkpoints/} contiene los
cinco \textit{checkpoints} de validación cruzada
(\texttt{gnn\_fold\_0.pt}--\texttt{gnn\_fold\_4.pt}), el \textit{checkpoint}
\texttt{gnn\_best.pt}, los dos modelos \ac{RF} serializados en formato
\texttt{.skops} (\texttt{rf\_model.skops} y \texttt{rf\_severity.skops}), los
pesos del \textit{ensemble} en \texttt{ga\_weights.json} y los ficheros de
métricas (\texttt{cv\_results.json}, \texttt{rf\_results.json},
\texttt{ga\_results.json} y \texttt{test\_results.json}). El conjunto ocupa
aproximadamente 32~MB y es el mismo que consume la \ac{CLI} a través del enlace
simbólico \texttt{kubescan/checkpoints/trained/}.

Para garantizar la reproducibilidad, los resultados finales de esta memoria se
entrenaron y evaluaron sobre una GPU NVIDIA~T4 (CUDA) con Python~3.12, PyTorch~2.11
(CUDA~12.8), PyTorch Geometric~2.8, scikit-learn~1.6 y NumPy~2.0. Esa
configuración corresponde al entorno de entrenamiento e investigación, que exige
GPU y las dependencias completas de \textit{PyTorch Geometric}; el paquete
distribuible admite un rango de instalación más amplio y su instalación se
verificó también en entornos Python~3.10 y~3.11, según se detalla en la
Sección~\ref{sec:verificacion_requisitos} (RNF-3). Las versiones
exactas de todas las dependencias, junto con las semillas, las particiones y las sumas
de verificación (\textit{checksums}) de los datos y los \textit{checkpoints}, se
registran en \texttt{research/models/checkpoints/run\_manifest.json}; el flujo completo
se reproduce mediante los objetivos del \texttt{Makefile} (\texttt{make data},
\texttt{make reproduce}).

El estudiante es el único autor de todos los \textit{commits} del repositorio.
Los datos públicos utilizados provienen de las fuentes citadas en la
Sección~\ref{sec:fuentes_datos} (Rahman et~al., BishopFox BadPods, Kubernetes Goat y repositorios
de seguridad adicionales de dominio público).
